% =====================================================================
% Problema aplicado --- Ravi (Frente 5)
% Só o Ravi edita este arquivo.
% =====================================================================
\section{Problema aplicado}\label{sec:problema}

Queremos saber quais streamers da Twitch PT são os influenciadores e qual centralidade os encontra
melhor. A pergunta se divide em duas: a \textbf{P1} pergunta qual centralidade melhor ordena os
streamers pelo alcance de um espalhamento simulado que começa em cada um, e a \textbf{P2}, qual melhor
os ordena pelas visualizações reais do canal. A hipótese \textbf{H1} é que o número de núcleo supera
o grau na P1, como já se observou em outras redes~\cite{kitsak2010}. A hipótese \textbf{H2} é que
todas as centralidades concordam com as visualizações, mas menos do que com o alcance, porque a
audiência também depende de fatores que o grafo não registra, como a idade do canal.

\paragraph{Entrada e saída.} A entrada é o grafo $G = (V, E)$ da Twitch PT (Seção~\ref{sec:bases}),
com $n = 1.912$ streamers e $m = 31.299$ amizades mútuas, e as visualizações $w(v)$ de cada canal. São
comparadas seis centralidades $c$: grau, proximidade e harmônica (Seção~\ref{sec:caminhos}),
intermediação, PageRank e número de núcleo $k_s$ (Seção~\ref{sec:centralidades}). A saída é o
ranking dos streamers pelo alcance $\rho$, com os empates desfeitos pelo grau e depois pelo
identificador do nó, e, para cada $c$, dois números: $\tau_b(c, \rho)$, a concordância do ranking de
$c$ com o alcance, que responde à P1, e $\tau_b(c, w)$, a concordância com as visualizações, que
responde à P2.

\paragraph{Como validar.} O alcance $\rho(v)$ é a fração média da rede atingida por uma simulação SIR
em tempo discreto que começa só em $v$, repetida $M = 100$ vezes: a cada passo, cada infectado
transmite para cada vizinho suscetível com probabilidade $\beta$ e se recupera no passo seguinte
($\gamma = 1$). Abaixo do limiar epidêmico
$\beta_c = \langle k \rangle / (\langle k^2 \rangle - \langle k \rangle) \approx 0{,}0079$~\cite{newman2002},
todo surto morre pequeno, e muito acima dele quase toda a rede é atingida; nos dois casos, o alcance
não distingue os nós. Por isso, usamos $\beta = 0{,}012 \approx 1{,}5\,\beta_c$, pouco acima do
limiar, como em~\cite{kitsak2010}. A concordância entre uma centralidade $c$ e uma referência
$Y \in \{\rho, w\}$ é o $\tau_b$ de Kendall, a variante que trata os empates~\cite{kendall1945}:
\[
  \tau_b(c, Y) = \frac{P - Q}{\sqrt{(P + Q + T_c)(P + Q + T_Y)}},
\]
em que, entre os $\binom{n}{2}$ pares de nós, $P$ e $Q$ contam os pares concordantes e discordantes,
e $T_c$ e $T_Y$, os empatados só em $c$ ou só em $Y$. O $\tau_b$ só olha a ordem, o que permite
comparar medidas em escalas diferentes, e tratar os empates é essencial, porque o grau assume só 187
valores distintos e $k_s$, só 37. Assim, \textbf{H1}: $\tau_b(k_s, \rho) > \tau_b(\text{grau}, \rho)$;
e \textbf{H2}: $0 < \tau_b(c, w) < \tau_b(c, \rho)$ para toda $c$.
